% F5 of article/paper-figures.md: the flow-first construction, one period,
% top to bottom.  The object passing between the steps is drawn, not named,
% and every band carries the same two-line statement of what it fixes and
% what it leaves free.  Greyscale only (fills black!X plus one hatch), so
% nothing depends on colour.  The libraries used -- patterns,
% decorations.pathreplacing, arrows.meta -- are all loaded by
% ../necessarypackages and paper.tex; no new package is needed.
\begin{figure}[!htb]
\centering
\resizebox{\textwidth}{!}{%
\begin{tikzpicture}[>=Latex, line width=0.5pt,
  band/.style={draw=black!45, fill=black!2, rounded corners=2pt},
  ttl/.style={font=\small\bfseries, anchor=west, align=left},
  bdg/.style={circle, draw=black!70, fill=white, inner sep=0pt,
              minimum size=4.6mm, font=\scriptsize\bfseries},
  ann/.style={font=\scriptsize, align=center},
  annl/.style={font=\scriptsize, align=left},
  mic/.style={font=\tiny, align=center},
  micl/.style={font=\tiny, align=left},
  st/.style={font=\scriptsize, align=left, anchor=west, text width=38mm},
  barO/.style={draw=black!70, fill=white},         % still continuous
  qty/.style={draw=black!70, fill=black!25},
  trk/.style={draw=black!80, fill=white, rounded corners=0.6pt},
  seg1/.style={draw=black!55, fill=black!30},      % R1 units
  seg2/.style={draw=black!55, fill=black!10},      % R2 units
  sur/.style={draw=black!65, pattern=north east lines, pattern color=black!75},
  glyA/.style={draw=black!55, fill=black!12},      % a column already in the pool
  glyB/.style={draw=black!70, fill=black!35},      % a part built here
  rule/.style={draw=black!35, dotted}]

\def\BW{15.2}          % content width
\def\SX{10.95}         % the fixed/free rule
\def\GX{2.20}          % x of the between-band arrows

% =====================================================================
% band 1 -- the node's LP
% =====================================================================
\draw[band] (0,0) rectangle (\BW,-2.25);
\draw[rule] (\SX,-0.12) -- (\SX,-2.13);
\node[bdg] at (0.55,-0.38) {1};
\node[ttl] at (0.95,-0.38) {Restricted master, inside $\mathcal{R}$};

\draw[black!55] (0.60,-1.82) -- (4.55,-1.82);
\foreach \i/\h/\v in {0/0.95/{1.0}, 1/0.86/{0.9}, 2/0.67/{0.7},
                      3/0.48/{0.5}, 4/0.29/{0.3}} {
  \pgfmathsetmacro{\x}{0.85+\i*0.72}
  \draw[barO] (\x,-1.82) rectangle (\x+0.46,-1.82+\h);
  \node[mic, anchor=south] at (\x+0.23,-1.79+\h) {\v};
}
\node[mic] at (2.58,-2.00) {one bar per column $\omega\in\bar\Omega$};
\node[annl, anchor=west] at (4.80,-1.38)
  {$\sum_{\omega}\lambda^{g}_{\omega t}=3.4$};
\node[mic, anchor=north west] at (4.80,-1.62) {not an integer};

\node[ann] at (7.55,-0.98) {duals $\pi_{ct},\,\mu_{gt}$};
\draw[->] (7.55,-1.20) -- (7.55,-1.76);
\node[mic, anchor=north] at (7.55,-1.80) {to the pricing problem};

\node[ann, text width=22mm] at (9.60,-1.26)
  {on a \emph{fresh copy} of the RMP: step~2 changes variable types};

\node[st] at (11.15,-1.12)
  {\textbf{fixed:} the region $\mathcal{R}$ --- bounds on $n_{gt}$ and $X$\\[2pt]
   \textbf{free:} every $\lambda^{g}_{\omega t}$, $\overline{V}_{nt}$,
   $\overline{V^{\mathrm{e}}}_{rt}$, all continuous};

% ---- 1 -> 2 ----------------------------------------------------------
\draw[->, line width=0.8pt] (\GX,-2.32) -- (\GX,-2.80);
\node[micl, anchor=west] at (\GX+0.22,-2.56)
  {the same rows; only the variable types change};

% =====================================================================
% band 2 -- the aggregate integer master
% =====================================================================
\draw[band] (0,-2.87) rectangle (\BW,-5.42);
\draw[rule] (\SX,-2.99) -- (\SX,-5.30);
\node[bdg] at (0.55,-3.25) {2};
\node[ttl] at (0.95,-3.25) {Aggregate integer master};

\draw[black!55] (0.60,-5.00) -- (4.55,-5.00);
\foreach \i/\h/\v in {0/0.90/{1.2}, 1/0.68/{0.9}, 2/0.60/{0.8},
                      3/0.53/{0.7}, 4/0.30/{0.4}} {
  \pgfmathsetmacro{\x}{0.85+\i*0.72}
  \draw[barO] (\x,-5.00) rectangle (\x+0.46,-5.00+\h);
  \node[mic, anchor=south] at (\x+0.23,-4.97+\h) {\v};
}
\node[ann, anchor=south] at (2.58,-3.92)
  {$\sum_{\omega}\lambda^{g}_{\omega t}=n^{\star}_{gt}=4$};
\node[mic] at (2.58,-5.18) {the $\lambda$ themselves stay continuous};

\node[ann, anchor=south] at (6.60,-3.92) {$Q^{\star}_{gt}$, integer};
\node[mic, anchor=east] at (5.30,-4.33) {$R_{1}$};
\draw[qty] (5.38,-4.46) rectangle (7.38,-4.20);
\node[mic, anchor=west] at (7.46,-4.33) {200};
\node[mic, anchor=east] at (5.30,-4.83) {$R_{2}$};
\draw[qty] (5.38,-4.96) rectangle (5.68,-4.70);
\node[mic, anchor=west] at (5.76,-4.83) {30};
\node[mic] at (6.60,-5.18) {bar length $\propto$ units};

\node[ann, text width=25mm] at (9.48,-4.42)
  {\emph{no pricing call here}: the rows are the master's own, and the
   objective is decided now};

\node[st] at (11.15,-4.15)
  {\textbf{fixed:} $n^{\star}_{gt}$, and the integer deliveries
   $Q^{\star}_{gt}=(\overline{V}_{nt},\overline{V^{\mathrm{e}}}_{rt})$\\[2pt]
   \textbf{free:} the columns $\lambda^{g}_{\omega t}$, \emph{still} continuous};

% ---- 2 -> 3 ----------------------------------------------------------
\draw[->, line width=0.8pt] (\GX,-5.49) -- (\GX,-5.97);
\node[micl, anchor=west] at (\GX+0.22,-5.73)
  {$(Q^{\star}_{gt},\,n^{\star}_{gt})$ --- fixed over continuous $\lambda$,
   so \emph{not} a sum of whole legal loads};

% =====================================================================
% band 3 -- the partition, as a cover
% =====================================================================
\draw[band] (0,-6.04) rectangle (\BW,-11.46);
\draw[rule] (\SX,-6.16) -- (\SX,-8.98);
\node[bdg] at (0.55,-6.42) {3};
\node[ttl] at (0.95,-6.42) {Partition --- a cover, not a split};

% what has to be delivered
\draw[seg1] (0.72,-8.52) rectangle (1.32,-7.07);
\draw[seg2] (0.72,-7.07) rectangle (1.32,-6.77);
\node[mic] at (1.02,-7.80) {$R_{1}$};
\node[mic, anchor=west] at (1.40,-6.92) {$R_{2}$};
\node[mic, anchor=north] at (1.02,-8.58) {$Q^{\star}_{gt}$};
\draw[->] (1.95,-7.80) -- (2.62,-7.80);
\node[mic, anchor=east] at (2.56,-7.20) {$\alpha\Gamma_{K}$};
\node[mic, anchor=north] at (2.28,-7.88) {$n^{\star}_{gt}=4$};

% the four parts
\foreach \k/\f/\r/\s/\z/\b in {%
   1/1.56/1.12/0.14/{$z_{1}{=}1$}/{full by vol.},
   2/1.14/0.86/0.12/{$z_{2}{=}0$}/{full by wgt.},
   3/1.48/1.04/0.22/{$z_{3}{=}1$}/{full by vol.},
   4/1.06/0.78/0.14/{$z_{4}{=}0$}/{full by wgt.}} {
  \pgfmathsetmacro{\x}{2.75+(\k-1)*1.60}
  \pgfmathsetmacro{\yr}{-8.52+\r}
  \pgfmathsetmacro{\yf}{-8.52+\f}
  \pgfmathsetmacro{\ys}{\yf-\s}
  \draw[trk] (\x,-8.52) rectangle (\x+1.30,-6.77);
  \draw[seg1] (\x+0.06,-8.46) rectangle (\x+1.24,\yr);
  \draw[seg2] (\x+0.06,\yr) rectangle (\x+1.24,\ys);
  \draw[sur]  (\x+0.06,\ys) rectangle (\x+1.24,\yf);
  \draw[dashed, black!85] (\x-0.10,-7.20) -- (\x+1.40,-7.20);
  \node[mic, text width=15mm, anchor=north] at (\x+0.65,-8.60) {\z\\ \b};
}

% key
\draw[sur]  (9.14,-7.14) rectangle (9.44,-6.90);
\node[mic, anchor=west] at (9.52,-7.02) {surplus};
\draw[seg2] (9.14,-7.64) rectangle (9.44,-7.40);
\node[mic, anchor=west] at (9.52,-7.52) {$R_{2}$ units};
\draw[seg1] (9.14,-8.14) rectangle (9.44,-7.90);
\node[mic, anchor=west] at (9.52,-8.02) {$R_{1}$ units};

\node[st] at (11.15,-7.50)
  {\textbf{fixed:} $n^{\star}_{gt}$ trucks ($+1$ or $+2$ if none can be
   formed), and $Q^{\star}_{gtc}$ as a \emph{lower} bound\\[2pt]
   \textbf{free:} the counts $a_{kc}\le\bar u_{gc}$, and which resource
   binds ($z_{k}$)};

\node[ann, text width=144mm, anchor=north] at (7.45,-9.06)
  {$\textstyle\sum_{k}a_{kc}\ge Q^{\star}_{gtc}$ \eqref{eq:part-cover}:
   a \textbf{cover}, not a split. Each part is a legal full truck
   \eqref{eq:part-cap}--\eqref{eq:part-floor}, full on volume \textbf{or}
   on payload and never on the sum of the two; the surplus is minimised in
   volume \eqref{eq:part-obj}, and the parts are ordered by weight
   \eqref{eq:part-sym}.};

\node[ann, text width=139mm, draw=black!45, dashed, fill=black!5,
      rounded corners=1.5pt, inner sep=4pt, anchor=north] at (7.60,-10.18)
  {\emph{Why the inequality cannot be made an equality:}
   $\{R_{1}{:}\,200,\,R_{2}{:}\,30\}$ asks for $35.5\,\mathrm{m}^{2}$ of
   stacks on a $32.64\,\mathrm{m}^{2}$ floor, so not one truck, and holds
   $1.83$ full-truckload floors of volume, so not two full ones --- no
   number of full trucks splits it exactly.};

% ---- 3 -> 4 ----------------------------------------------------------
\draw[->, line width=0.8pt] (\GX,-11.53) -- (\GX,-12.01);
\node[micl, anchor=west] at (\GX+0.22,-11.77)
  {$a_{1},\dots,a_{n^{\star}_{gt}}$ --- each a legal full truck, hence a
   legal column};

% =====================================================================
% band 4 -- the integer master over pool and parts
% =====================================================================
\begin{scope}[shift={(0,-0.66)}]
\draw[band] (0,-11.42) rectangle (\BW,-14.30);
\draw[rule] (\SX,-11.54) -- (\SX,-14.18);
\node[bdg] at (0.55,-11.80) {4};
\node[ttl] at (0.95,-11.80) {Integer master over pool and parts};

\draw[draw=black!40, fill=black!3, rounded corners=1.5pt]
   (0.62,-12.34) rectangle (4.86,-13.06);
\foreach \i in {0,...,5} {
  \pgfmathsetmacro{\x}{0.78+\i*0.40}
  \draw[glyA] (\x,-12.90) rectangle (\x+0.32,-12.50);
}
\foreach \i in {6,...,9} {
  \pgfmathsetmacro{\x}{0.78+\i*0.40}
  \draw[glyB] (\x,-12.90) rectangle (\x+0.32,-12.50);
}
\node[ann, anchor=north] at (1.60,-13.12) {pool $\bar\Omega\,\cup\,\{a_{k}\}$};
\draw[decorate, decoration={brace, amplitude=3pt, mirror}, black!60]
   (3.18,-13.14) -- (4.70,-13.14);
\node[mic, anchor=north] at (3.94,-13.24) {the four parts};

\draw[->, line width=0.8pt] (5.02,-12.70) -- (6.52,-12.70);
\node[mic, text width=22mm, anchor=north] at (6.02,-12.90)
  {integer master, warm-started with the partition};

\foreach \k/\f/\l in {1/{glyB}/{$a_{1}$}, 2/{glyB}/{$a_{2}$},
                      3/{glyB}/{$a_{3}$}, 4/{glyA}/{$\omega\in\bar\Omega$}} {
  \pgfmathsetmacro{\x}{6.85+(\k-1)*1.00}
  \draw[\f] (\x,-13.06) rectangle (\x+0.75,-12.34);
  \node[mic] at (\x+0.375,-12.70) {\l};
}
\node[ann, anchor=south] at (8.73,-12.30) {the node's plan};

\draw[draw=black!45, dashed] (7.05,-14.12) rectangle (7.80,-13.62);
\node[mic] at (7.42,-13.87) {$a_{4}$};
\node[annl, text width=26mm, anchor=west] at (7.96,-13.87)
  {built, not dispatched};

\node[st] at (11.15,-12.90)
  {\textbf{fixed:} the pool $\bar\Omega\cup\{a_{k}\}$, and the region
   $\mathcal{R}$\\[2pt]
   \textbf{free:} integer $\lambda^{g}_{\omega t}$ --- how many of each
   column goes out};
\end{scope}

\end{tikzpicture}}
\caption{The flow-first construction over one period, and why its third
step covers rather than splits. Steps~1 and~2 settle how much moves and
how many trucks carry it while the column variables stay
\emph{continuous}, so the delivery vector they fix need not be a sum of
whole legal loads --- for the vector shown, no number of full trucks
realises it exactly. Step~3 therefore asks only that the parts
\emph{cover} it \eqref{eq:part-cover}, paying for the hatched surplus
\eqref{eq:part-obj}, and makes each part full on volume \emph{or} on
payload \eqref{eq:part-floor} rather than on the sum of the two. Nothing
built is binding: the parts enter step~4 as columns with a warm start,
and the integer master may dispatch fewer of them than the cover
produced.}
\label{fig:construct}
\end{figure}
